% ============================================================
% 项目全景论证 — 问题 · 需求 · 技术 · 价值
% 面向客户/评审的完整全貌文档，一页看懂 + 逐层展开
% 编译: xelatex -interaction=nonstopmode 全景-问题需求技术价值.tex
% 口径对齐: 报价书(报价-软件定制开发.tex) / 附件② / 附件④ / 报价策略要点-V8.md
% ============================================================
\documentclass[11pt,a4paper]{ctexart}
\usepackage{xeCJK}
\usepackage{graphicx}
\usepackage{float}
\usepackage{fancyhdr}
\usepackage{hyperref}
\usepackage{geometry}
\usepackage{longtable}
\usepackage{booktabs}
\usepackage{enumitem}
\usepackage{caption}
\usepackage{array}
\usepackage{tikz}
\usetikzlibrary{arrows.meta, backgrounds}

% 字体
\setCJKmainfont{Microsoft YaHei}

% 页面
\geometry{a4paper, margin=2.5cm}

% 页眉页脚
\pagestyle{fancy}
\fancyhead[L]{时序数据采集与应用管理系统 — 项目全景论证}
\fancyhead[R]{问题 · 需求 · 技术 · 价值}
\fancyfoot[C]{\thepage}

\hypersetup{
  colorlinks=true, linkcolor=black, urlcolor=blue,
  pdftitle={项目全景论证——问题·需求·技术·价值},
}

\begin{document}

\begin{center}
{\Huge\bfseries 项目全景论证}\\[10pt]
{\Large 问题 · 需求 · 技术 · 价值 —— 一页看懂本项目}\\[8pt]
{\large 时序数据采集与应用管理系统 \quad 版本 V1.0 \quad 编制日期 2026-08}
\end{center}

\vspace{1em}
\noindent\hrulefill
\vspace{1em}

% ═══════════════════════════════════════════
\section{全景总览：一页看懂}

\subsection{一句话定位}

本系统是既有 A11 油气生产物联网系统的\textbf{能力增强与扩展}：\textbf{不改 A11、不替代 A11}，以\textbf{路由监听双路采集}方式，在保留 A11 原有业务完全不变的前提下，为油田提供\textbf{高频、动态、差异化、信创化}的时序数据采集与应用能力——\textbf{全域感知、多维联动}。

\subsection{四问因果链}

本项目的一切论证围绕四个问题展开，且\textbf{四问环环相扣}：问题的定性决定需求清单，需求清单约束技术路径，技术路径兑现价值点。

\begin{figure}[H]\centering
\begin{tikzpicture}[
  font=\small,
  box/.style={draw, rounded corners=2pt, align=center, inner sep=4pt},
  layer/.style={box, fill=gray!8, font=\bfseries},
  item/.style={box, fill=white, text width=4.6cm, font=\scriptsize},
  arrow/.style={-{Stealth}, very thick, black!70},
  lab/.style={font=\footnotesize\bfseries, black!70}
]
% ═══ ① 问题 ═══
\node[layer, text width=15.0cm] (p1) at (7.8,14.2) {① 问题 —— 为什么做};
\node[item] (p1a) at (2.6,12.9) {\textbf{核心矛盾}\\油田精细化/定制化需求 vs\\A11 标准固化体系\\（统一点位·统一频率·一种节奏）};
\node[item] (p1b) at (7.8,12.9) {\textbf{双重保留挑战}\\A11 保留运行、不替代\\全部既有配置与业务集成\\原样保留、不中断、不重建};
\node[item] (p1c) at (13.0,12.9) {\textbf{硬约束}\\新增硬件仅 2--3 台\\其余设施全部利旧\\IO 服务器共存部署};
\draw[arrow] (7.8,11.9) -- (7.8,10.8);
\node[lab] at (8.4,11.35) {问题定需求};

% ═══ ② 需求 ═══
\node[layer, text width=15.0cm] (p2) at (7.8,10.4) {② 需求 —— 前置约束 15 条 + 操作纪律 4 条（全部不可妥协）};
\node[item, text width=3.5cm] (p2a) at (1.9,9.1) {\textbf{系统保留（5 条）}\\A11 保留运行·不改 DTU\\既有配置保留·集成不断\\原有业务零影响};
\node[item, text width=3.5cm] (p2b) at (5.7,9.1) {\textbf{评审硬性（4 条）}\\路由监听双路采集\\主数据唯一源头\\两层部署·源码+外协限三类};
\node[item, text width=3.5cm] (p2c) at (9.5,9.1) {\textbf{资源商务（4 条）}\\硬件 2--3 台·全利旧\\共存轻量（CPU $<$5\%）\\商务口径与官方测算一致};
\node[item, text width=3.5cm] (p2d) at (13.3,9.1) {\textbf{合规手续（2 条）}\\甲方书面授权委托\\生产网纪律（只读·禁装·禁停）};
\draw[arrow] (7.8,8.1) -- (7.8,7.0);
\node[lab] at (8.4,7.55) {需求定技术};

% ═══ ③ 技术 ═══
\node[layer, text width=15.0cm] (p3) at (7.8,6.6) {③ 技术 —— 路由监听双路采集（约束的必然结果）};
\node[item] (p3a) at (2.6,5.3) {\textbf{静态 IP（设备侧 server）}\\客户端主动连接采集（旁路）\\不改设备·设备无感知\\原 A11 链路原样不动};
\node[item] (p3b) at (7.8,5.3) {\textbf{动态 IP（设备侧 client）}\\抢占原 IP:端口桥接\\复制一份+原样转发\\唯此一路径全覆盖};
\node[item] (p3c) at (13.0,5.3) {\textbf{全程安全}\\三阶段渐进·低峰窗口秒级切换\\业务流摸清·异常即时回退\\零影响可验收};
\node[box, fill=green!8, text width=15.0cm] (p3d) at (7.8,3.8) {统一汇聚 $\rightarrow$ 边缘流式计算 $\rightarrow$ DMZ 短期缓存 $\rightarrow$ \textbf{数据湖（主数据唯一源头）}};
\draw[arrow] (7.8,2.8) -- (7.8,1.9);
\node[lab] at (8.4,2.35) {技术兑现价值};

% ═══ ④ 价值 ═══
\node[layer, text width=15.0cm] (p4) at (7.8,1.5) {④ 价值 —— 全域感知 · 多维联动（A11 管存量，本系统管增量）};
\node[item, text width=3.5cm] (p4a) at (2.0,0.2) {\textbf{高频}\\500ms--60s 动态调频\\异常加密 10min$\rightarrow$秒级};
\node[item, text width=3.5cm] (p4b) at (5.9,0.2) {\textbf{动态}\\10 家网关厂家适配\\动态感知·断网缓存补传};
\node[item, text width=3.5cm] (p4c) at (9.8,0.2) {\textbf{差异化}\\15 种算法·逐井调优\\告警入库前闭环};
\node[item, text width=3.5cm] (p4d) at (13.7,0.2) {\textbf{信创+安全}\\麒麟/飞腾全栈适配\\手术级安全·全栈源码交付};
\end{tikzpicture}
\caption{四问全景图——问题 $\rightarrow$ 需求 $\rightarrow$ 技术 $\rightarrow$ 价值，环环相扣：问题定需求，需求定技术，技术兑现价值}
\end{figure}

\subsection{推导逻辑（全貌的"灵魂"）}

\begin{enumerate}[leftmargin=*]
  \item \textbf{问题定需求}：矛盾与双重保留，逐条推导出 15 条前置约束（第 3 节）；
  \item \textbf{需求定技术}：15 条约束中任何一条都不允许"改 A11 / 改 DTU / 换设备"，技术上只剩唯一路径——\textbf{旁路复制（桥接/镜像）}，没有任何技术偏好，是约束的必然结果；
  \item \textbf{技术兑现价值}：旁路复制带来的高频数据密度、动态感知、差异化调优能力，正是 A11 给不了的、油田要的东西（第 5 节）。
\end{enumerate}

\textbf{一句话全貌}：\textbf{A11 动不得 ⟹ 只能旁路拿数据 ⟹ 拿到了 A11 拿不到的数据 ⟹ 做 A11 做不了的应用。}

% ═══════════════════════════════════════════
\section{问题：我们面对的是什么}

\subsection{核心矛盾}

油田近年全面推进电参数字化、工况智能诊断、设备健康预警等应用，运行管理日益\textbf{精细化、定制化}；而 A11 是集团统建系统，体系\textbf{标准固化}——统一点位、统一频率、统一配置，全系统一种节奏：

\begin{itemize}[leftmargin=*]
  \item 数据分钟级定时抽吸（约 10--30 分钟一档），功图类曲线约 30 分钟/副——高频异常（皮带断裂、电潜泵参数突变、频率跃升冲击）在数据刷新间隙内\textbf{无法被识别}；
  \item 报警依赖后台批量处理，无本地流式计算，\textbf{无法做到秒级/亚秒级异常响应}；
  \item 差异化报警（逐井、逐工况）\textbf{无逐井调优通道}；
  \item 动态接入（移动网关、设备增减）\textbf{无统一感知机制}；
  \item 协议支撑有限，视频流、智能仪表、新能源等新兴设备源\textbf{无法接入}；
  \item 依赖老旧 Windows 环境，\textbf{无国产 CPU/OS 适配能力}。
\end{itemize}

\textbf{瓶颈不在算法，而在数据密度与采集灵活性}——算法能力在局地试点已获验证，受限于分钟级采集底座无法落地。

\subsection{双重保留挑战（本项目最难的约束）}

在解决矛盾的同时，还必须满足两条"保留"要求，任何一条做不到，方案即不成立：

\begin{enumerate}[leftmargin=*]
  \item \textbf{保留 A11 运行}：A11 继续保留运行、继续承载原有生产业务，不替代、不改动、不中断；
  \item \textbf{保留现有系统全部既有配置与业务集成}：已建点位、量程频次、物模型、A2/A5 资产体系全部原样保留，与既有业务系统的数据交换、接口对接关系不中断、不重建。
\end{enumerate}

\subsection{硬约束：投入极少、全部利旧}

\begin{itemize}[leftmargin=*]
  \item 项目新增硬件仅 \textbf{2--3 台服务器}（生产网边缘计算 1 台、DMZ 短期缓存 1 台、办公网汇聚中枢 1 台，可按现场合并为 2 台）；
  \item 其余设施——IO 服务器、网络、存储、既有设备——\textbf{全部利旧}；
  \item 本系统软件\textbf{直接部署于既有 IO 服务器}，与 A11 服务\textbf{同机共存、互不影响}（轻量只读采集组件，CPU 占用 $<$5\%、延迟增加 $<$20$\mu$s，独立进程、独立端口、资源限额、故障自动静默）。
\end{itemize}

% ═══════════════════════════════════════════
\section{需求：精准需求（15 条前置约束 + 4 条操作纪律）}

\subsection{约束总表}

经系统梳理，本项目前置约束共 \textbf{15 条，分 4 组}，全部不可妥协；另由桥接抢占形态派生 \textbf{4 条操作纪律需求}（见 3.3）：

\begin{table}[H]\centering
\caption{前置约束总表（15 条 / 4 组）}
\begin{tabular}{p{2.4cm} p{2.2cm} p{9.2cm}}
\toprule
\textbf{组别} & \textbf{条数} & \textbf{约束内容} \\
\midrule
一、系统保留 & 5 & A11 保留运行；不改 DTU 配置；保留全部既有配置（点位/量程频次/物模型/A2--A5 资产体系）；业务系统集成关系不中断；原有业务零影响 \\
二、评审硬性 & 4 & 路由监听双路采集、不改 A11；主数据唯一源头；两层部署（DMZ 只转发不存）；全栈源码交付 + 外协限三类 \\
三、资源商务 & 4 & 硬件仅 2--3 台；其余全利旧；IO 服务器共存轻量部署（CPU $<$5\%、延迟 $<$20$\mu$s）；商务与官方功能点测算一致 \\
四、合规手续 & 2 & 甲方书面授权委托；生产网操作纪律（只读观测、禁止安装/重启/停服） \\
\bottomrule
\end{tabular}
\end{table}

\subsection{约束链：每条约束 ⟹ 方案一个条款}

\textbf{本项目的技术方案不是"设计出来的偏好"，而是"15 条约束推导出来的必然结果"}——每条约束都能在方案中找到对应条款：

\begin{longtable}{p{1.2cm} p{4.8cm} p{7.8cm}}
\toprule
\textbf{约束} & \textbf{内容} & \textbf{方案对应条款} \\
\midrule
\endfirsthead
\multicolumn{3}{c}{(续表)} \\
\toprule
\textbf{约束} & \textbf{内容} & \textbf{方案对应条款} \\
\midrule
\endhead

\multicolumn{3}{l}{\textbf{组一 系统保留（约束 1--5）}} \\
1 & A11 保留运行 & 双路采集：原 A11 链路原样运行，本系统旁路采集同一数据流，实时同步（附件② 4.1） \\
2 & 不改 DTU 配置 & 设备侧数据流只能\textbf{桥接复制}获取，全程不存在"接管"（附件② 4.1 接入节奏） \\
3 & 保留全部既有配置 & A11 侧零配置变更，本系统以主数据映射方式对接既有点位（附件② 第 5 节） \\
4 & 业务集成不中断 & 数据湖门户嵌入方式融入业务系统生态，业务系统零感知（附件② 第 1 节） \\
5 & 原有业务零影响 & 并联无阻断、桥接仅原样续传、故障静默；三阶段渐进可回退；"实施前后 A11 运行对比记录"列入验收（报价书 7.2） \\
\midrule
\multicolumn{3}{l}{\textbf{组二 评审硬性（约束 6--9）}} \\
6 & 路由监听双路采集不改 A11 & 静态 IP $\rightarrow$ 客户端主动采集（旁路）；动态 IP $\rightarrow$ 抢占原 IP:端口桥接（附件② 4.1） \\
7 & 主数据唯一源头 & 点位/物模型全部注册至数据湖主数据体系，不建第二套体系（附件② 第 5 节） \\
8 & 两层部署 & 生产网边缘采集不落地长周期数据，DMZ 只转发不存，办公网入湖（附件② 第 6 节） \\
9 & 全栈源码交付 + 外协限三类 & 自研 790 人天 / 外协 550 人天（限协议适配、安全认证、主数据对接）/ 开源复用 110 人天（附件④ 第 4 节） \\
\midrule
\multicolumn{3}{l}{\textbf{组三 资源商务（约束 10--13）}} \\
10 & 硬件仅 2--3 台 & 边缘计算/缓存/汇聚三节点，可合并为 2 台（附件② 4.1 硬件投入） \\
11 & 其余全利旧 & IO 服务器、网络、存储、既有设备全部利旧（报价书报价边界） \\
12 & IO 服务器共存部署 & 轻量采集组件：静态客户端采集 + 动态桥接转发，CPU $<$5\%，延迟 $<$20$\mu$s，故障自动静默（附件② 4.1） \\
13 & 商务与官方测算一致 & 5 板块 48 计价项，与官方功能点测算（4,088.57 FP）金额一致（报价书） \\
\midrule
\multicolumn{3}{l}{\textbf{组四 合规手续（约束 14--15）}} \\
14 & 甲方书面授权委托 & 合规边界：授权 + 合同范围 + 数据归属甲方 + 最小化留存（附件② 4.3 合规性论证） \\
15 & 生产网操作纪律 & 只读验证、禁止安装/重启/停服，现场操作全部经甲方书面批准（生产网规则） \\
\bottomrule
\end{longtable}

\subsection{操作纪律需求（N1--N4，由桥接抢占派生）}

桥接层启用（抢占原地址端口）需在生产系统上执行受控切换且保证无损，派生 4 条操作纪律需求——\textbf{从"技术可行性约束"升级为"运维可操作性约束"}：

\begin{table}[H]\centering
\caption{操作纪律需求（N1--N4）}
\begin{tabular}{p{1.0cm} p{3.0cm} p{7.0cm} p{2.8cm}}
\toprule
\textbf{编号} & \textbf{需求} & \textbf{内容} & \textbf{性质} \\
\midrule
N1 & 切换窗口极短 & 抢占切换窗口压缩至秒级（目标毫秒级），低峰窗口执行 & 性能指标级，列入验收 \\
N2 & 业务流全面摸底 & 设备/会话/端口/保活/流量特征逐项建档，切换前完成 & 独立工作项，现场调研前置 \\
N3 & 回退预案与演练 & 秒级切回原 IP:端口，全程值守，回退演练留痕 & 验收项（附件② 4.3） \\
N4 & 切换协同机制 & 切换方案甲方确认、窗口批准、三方值守（甲方运维、原 A11 服务方、我方） & 流程约束，升级约束 14 \\
\bottomrule
\end{tabular}
\end{table}

\textbf{需求定性变化}：动生产系统意味着责任判定需求显性化——甲方签署切换方案、我方出具过程记录，责任边界随审批流 + 记录流清晰化。

\subsection{约束落地：全部收敛为两大核心能力}

15 条约束不是分散的承诺，而是\textbf{全部收敛体现为两大产品化核心引擎}——\textbf{采集策略多样化}（怎么采）与\textbf{定制化流水计算}（怎么算）：

\begin{itemize}[leftmargin=*]
  \item \textbf{采集策略多样化（怎么采）}：双路采集按终端类型分流（静态客户端采集 / 动态桥接抢占）、三阶段渐进零影响接入、500ms--60s 可配置动态调频（异常加密 10min$\rightarrow$秒级 / 低峰降频 / 弱网退避）、10 家网关厂家适配、动态 IP 感知、断网缓存补传、按作业区独立定制；
  \item \textbf{定制化流水计算（怎么算）}：15 种内置算法 + 每作业区定制算法、入库前流水计算（告警入库前完成）、逐井差异化调优、物模型（G1--G8）驱动计算、边缘侧轻量执行（单条 $<$1ms、CPU $<$5\%）。
\end{itemize}

\begin{table}[H]\centering
\caption{四组约束 $\rightarrow$ 两大核心能力落地映射}
\begin{tabular}{p{2.2cm} p{5.8cm} p{5.8cm}}
\toprule
\textbf{约束组} & \textbf{落地于采集策略多样化} & \textbf{落地于定制化流水计算} \\
\midrule
系统保留 & 采集形态按终端分流；三阶段渐进接入，全程可回退 & 计算按主数据映射点位执行；旁路执行，不注入不干扰 \\
评审硬性 & 路由监听双路采集；生产网边缘采集、DMZ 只转发 & 主数据唯一源头驱动算法与告警；入库前流水计算 \\
资源商务 & 利旧共存轻量采集（CPU $<$5\%、延迟 $<$20$\mu$s） & 轻量计算引擎单条 $<$1ms、资源限额、故障静默 \\
合规手续 & 只读授权开发对接（凭证接入） & 只读计算、零注入、最小化留存 \\
\bottomrule
\end{tabular}
\end{table}

\textbf{客户视角一句话}：约束决定"怎么采、怎么算"——15 条约束逐条映射到两大引擎，两大引擎就是约束的\textbf{产品化载体}，承诺不落空、可验收。

% ═══════════════════════════════════════════
\section{技术：核心技术方案}

\subsection{总体架构：一条原链路 + 一条旁路}

\begin{itemize}[leftmargin=*]
  \item \textbf{原链路（不动）}：A11 采集链路原样运行，设备数据全量流向 A11，一切照旧；
  \item \textbf{旁路（新增）}：在同一数据流上并行复制一份镜像副本：采集 $\rightarrow$ 边缘流式计算 $\rightarrow$ DMZ 短期缓存 $\rightarrow$ 数据湖统一长期存储；
  \item \textbf{数据体系}：采集数据统一注册至数据湖\textbf{主数据体系}（唯一源头），不形成第二套体系；
  \item \textbf{跨网安全}：生产网只出不进，国密双向认证通道，控制指令穿透有独立审批。
\end{itemize}

\subsection{采集形态按终端接入类型分流}

\textbf{不是所有终端都能用同一种方式旁路}——按接入类型分流，两种形态并存、统一汇聚：

\begin{table}[H]\centering
\caption{采集形态分流}
\begin{tabular}{p{2.8cm} p{5.0cm} p{5.0cm}}
\toprule
\textbf{终端类型} & \textbf{形态} & \textbf{说明} \\
\midrule
静态 IP 终端（有线/固定链路，设备侧为服务端） & \textbf{旁路方案}：客户端主动连接采集 & 数据端点在设备侧（server），本系统以\textbf{客户端身份主动连接采集}：不改设备配置、设备侧无感知、原 A11 链路原样不动 \\
动态 IP 终端（4G/移动网关，设备侧为客户端） & \textbf{桥接方案}：抢占原服务器 IP 与端口 & 设备主动建连至原服务器，不改设备层 IP，\textbf{必须将原服务器 IP:端口抢占做桥接}：复制一份 + 透明转发原 A11 服务——唯一路径 \\
\bottomrule
\end{tabular}
\end{table}

\subsection{接入节奏：三阶段渐进，全程可回退}

\textbf{观察}（bypass，仅只读观察、不介入）$\rightarrow$ \textbf{桥接复制}（mirror，桥接层绑定原地址端口、复制并原样转发，原链路业务不断）$\rightarrow$ \textbf{稳定并行}（stable）——任一阶段异常\textbf{即时回退}，原链路零影响。

\subsection{对 A11 原有业务的零影响保障}

\begin{itemize}[leftmargin=*]
  \item \textbf{并联而非串联}：无阻断点，数据流始终全量流向 A11；
  \item \textbf{不修改任何配置}：A11 侧零配置变更；
  \item \textbf{不注入任何数据}：桥接仅原样续传原链路数据，不产生、不修改任何业务内容；
  \item \textbf{故障隔离与快速回退}：桥接层任何故障不影响业务续传——原 A11 服务秒级切回原 IP:端口，链路自动恢复；
  \item \textbf{资源可控}：独立进程、独立端口、资源限额，CPU $<$5\%、延迟 $<$20$\mu$s；
  \item \textbf{验收可证}：实施前后 A11 运行对比记录（采集成功率、链路时延与基线一致）作为验收项。
\end{itemize}

\subsection{桥接模式的技术难度（高难度要求）}

在"不改 DTU 配置、不改变任何端点收发行为"的硬约束下，\textbf{动态 IP 终端只能以桥接方式采集——桥接是本项目主形态而非边角场景}，这是本项目核心技术难点：

\begin{itemize}[leftmargin=*]
  \item \textbf{原 IP 与端口抢占}：桥接层接管原服务器 IP 与端口，原 A11 服务迁至备用地址/端口——切换与甲方协调、选低峰窗口执行，全程可回退；
  \item \textbf{透明应答与转发}：桥接层代替原服务完成 TCP 建连应答，数据复制一份后原样转发原 A11 服务，设备连接行为零感知；
  \item \textbf{TCP 流完整重组}：乱序、重传、分段、粘连按会话重组为完整应用层数据，重组错误即数据错位；
  \item \textbf{时序与顺序严格保持}：副本与源流一致，供流式计算与历史回放使用；
  \item \textbf{长连接会话跟踪}：设备重连、端口漂移、会话重建持续跟踪识别；
  \item \textbf{故障隔离与快速回退}：桥接层任何故障不影响业务续传——原服务秒级切回原 IP:端口，链路自动恢复。

\end{itemize}

对应工作量显式单列（附件④）：路由监听分析（A11 桥接链路）45 人天 + A11 专有协议帧解析 45 人天 + 试点联调 A11 桥接验证 50 人天。

\subsubsection{桥接切换操作纪律}

桥接层启用遵循\textbf{三条操作纪律}：\textbf{窗口期要非常短}（低峰窗口执行，压缩至秒级、目标毫秒级，业务可感知中断趋近于零）；\textbf{切换前业务流摸清}（设备清单、连接会话、端口占用、心跳保活、流量特征逐项建档，切换方案逐项核对、不留盲区）；\textbf{异常即时回退}（全程值守，任何异常按预案秒级切回原 IP:端口、链路自动恢复，出具切换前后对比记录验证无损）。

\subsection{合规性依据}

\begin{enumerate}[leftmargin=*]
  \item \textbf{合法性基础}：《网络安全法》第二十一条要求网络运营者"监测、记录网络运行状态"——采集生产数据以支撑安全生产监测，属法定职责；
  \item \textbf{国标认可的部署形态}：GB/T 20945--2023 附录 A 将审计产品"并联部署（网络数据镜像或分流）"与串联部署并列，为标准认可的旁路取数正规形态；
  \item \textbf{只读获取是行业标准做法}：NIST SP 800--82r3《OT 安全指南》（2023）明确"OT 环境中基于网络的监测能力通常经 SPAN 端口或被动网络 TAP 部署"；IEC 62443-3-3 SR 6.1 要求审计日志经授权只读访问——本项目获取通道只读、不写、不注入，与之一致；
  \item \textbf{等保义务}：GB/T 22239--2019（等保 2.0）要求在网络边界与重要节点实施安全审计——本项目以只读授权方式获取数据，同时也是协助履行等保审计义务；
  \item \textbf{合规边界}：甲方书面授权 + 合同范围约定 + 数据归属甲方 + 最小化留存、加密脱敏处置。
\end{enumerate}

% ═══════════════════════════════════════════
\section{价值：价值点}

\subsection{价值定位：A11 管存量，本系统管增量}

\textbf{核心能力总纲：全域感知、多维联动}——\textbf{全域感知}：全终端（静态/动态 IP 全覆盖）、全协议、全量高频采集，动态设备统一感知；\textbf{多维联动}：多源数据入库前联合判定、告警分级闭环通知、边云协同。

\textbf{一句话概括互补关系}：A11 管住存量（统建标准采集、既有生产业务不变），本系统管住增量（高频、动态、差异化、信创）——两条链路并行，数据汇聚于数据湖主数据体系。

\subsection{功能互补矩阵}

\begin{table}[H]\centering
\caption{与 A11 的功能互补对照}
\begin{tabular}{p{2.6cm} p{4.2cm} p{4.2cm} p{2.6cm}}
\toprule
\textbf{能力维度} & \textbf{A11 现状} & \textbf{本系统补充能力} & \textbf{对应模块} \\
\midrule
采集频率 & 分钟级定时抽吸 & 可配置 500ms--60s，异常自动加密（10min$\rightarrow$秒级）、低峰降频 & 定制动态采集 \\
边缘计算 & 无本地流式引擎，报警依赖后台批处理 & 入库前流水计算，15 种内置算法 + 每作业区定制算法，单条 $<$1ms & 流式计算引擎 \\
协议兼容 & 有限工业协议 & MQTT/\allowbreak{}HTTP/\allowbreak{}UDP/\allowbreak{}RTSP 多协议，10 家网关适配，非标协议按需接入 & 定制动态采集 \\
信创适配 & 依赖 Windows 老旧环境 & 全栈开源组件 + 自主定制，麒麟/飞腾等信创环境适配 & 全链路安全与审计追溯平台 \\
跨网贯通 & 生产网/办公网隔离 & 生产网边缘采集 + 办公网 DMZ 短期缓存入湖，安全通道穿透 & 实施交付 \\
\bottomrule
\end{tabular}
\end{table}

\subsection{高价值点 × 定制化对应（14 项）}

\textbf{每一个高价值收费点都有"开源/通用给不了、所以是定制"的论证}；通用能力公开承认开源复用、不计价：

\begin{longtable}{p{0.8cm} p{4.2cm} p{2.6cm} p{6.2cm}}
\toprule
\textbf{\#} & \textbf{高价值点} & \textbf{计价项} & \textbf{为什么是定制（开源给不了）} \\
\midrule
\endfirsthead
\multicolumn{4}{c}{(续表)} \\
\toprule
\textbf{\#} & \textbf{高价值点} & \textbf{计价项} & \textbf{为什么是定制（开源给不了）} \\
\midrule
\endhead

1 & A11 专有协议开发对接（psAPISDK/\allowbreak{}CommBridge 协议解析、Tag ID 发现） & 一-2、四-4.2 & 中石油私有协议，开源无实现，凭授权凭证开发对接 \\
2 & 路由监听双路采集（静态客户端采集/动态桥接转发、设备指纹/帧漂移检测） & 三-4、四-4.2 & 协议库开源，但双形态采集 + 协议识别 + 指纹漂移检测全自研 \\
3 & 入库前流水流式计算（15 种算法、单条 $<$1ms、告警入库前完成） & 一-5、一-6、四-4.3 & 开源流引擎重且非内存内判定；15 种油田算法集为自研资产 \\
4 & 智能动态调频（异常加密 10min$\rightarrow$秒级/低峰降频/弱网退避） & 四-4.1 & 开源采集框架无此油田业务策略，纯业务定制 \\
5 & 10 家网关厂家适配（注册帧识别/透传解析/心跳保活） & 四-4.1 & 各厂家私有透传协议，开源无，逐个适配开发 \\
6 & G1--G8 油水井物模型与自动配点（64 点/地址对齐） & 四-4.2 & 抽油机/螺杆泵/电潜泵物模型是油田行业特异资产，开源无 \\
7 & OPC DA 直采三级回退（DCS/PLC 差异、pSpace 比对 $<$1\%） & 四-4.2 & OPC DA 老技术，开源库不维护，现场环境恶劣需自研回退链 \\
8 & 国密算法套件（SM2/3/4 加密存储传输） & 一-9、二-5 & 信创合规刚需，通用开源无国密全栈，外协定制 \\
9 & 百万级测点高并发中枢（全油田 16 厂 · 100+ 作业区全部设备接入 · 260 万+ 点位满频采集） & 一-1、一-4 & 开源 MQTT 企业版才支撑此量级且收费；自主增强自研实现 \\
10 & 动态 IP 井下终端接入（4G 终端、端口映射穿透） & 二-4、四-4.1 & 油田 4G 井下场景特异，开源无现成支持 \\
11 & 物模型--主数据联动（注册制、唯一源头、自动配点） & 四-4.2、附件⑥ & 与数据湖主数据（A2/A5 资产体系）对接是油田特有，外协定制 \\
12 & 信创全栈适配（麒麟/UOS、飞腾/鲲鹏、达梦/金仓） & 五-4 & 社区版时序库不支持信创 OS，必须自主适配 \\
13 & 断网缓存与补传（时序补传/去重/幂等） & 四-4.1 & 油田弱网环境特异，开源无此业务闭环 \\
14 & 点位配置导出$\rightarrow$A11 导入（一套配置双向复用） & 四-4.1、四-4.2 & 评审点名功能，A11 格式特异，纯定制 \\
\bottomrule
\end{longtable}

\textbf{规律}：高价值点全部落在产品化封装（A11 专有协议/流式计算/告警引擎/国密/高并发）与作业区定制；部署与标准化组件是低价值复用点，敢公开承认——\textbf{没有一个是"纯通用开源"}。

\subsection{投入产出}

\begin{itemize}[leftmargin=*]
  \item \textbf{投入极小}：硬件仅 2--3 台服务器，其余全利旧；软件部署于既有 IO 服务器，无新建机房、无大规模设备采购；
  \item \textbf{工作量可审计}：1,450 人天 / 18--20 人团队 / 4 个月交付（2026-08 至 2026-11-30），与官方功能点测算口径一致；
  \item \textbf{与替代路径的对比}：若通过改造 A11 统建系统实现同等能力，投入将达亿元级（详见论证报告第三方证据链）；本项目以不足其零头的投入获得同等数据能力，\textbf{投入产出比跨数量级占优}；
  \item \textbf{资产可复用}：核心引擎（协议解析/流式计算/告警引擎）一次开发、多项目复用，占投资近四成；后续新作业区仅按"区数 $\times$ 单价"扩展，推广至其他油田时核心资产零成本复用；
  \item \textbf{商务金额}：详见《软件定制开发报价书》（与官方功能点测算 4,088.57 FP 折算金额一致）。
\end{itemize}

% ═══════════════════════════════════════════
\section{结论：四问收口}

\begin{table}[H]\centering
\caption{四问收口}
\begin{tabular}{p{1.6cm} p{12.2cm}}
\toprule
\textbf{四问} & \textbf{收口结论} \\
\midrule
问题 & 矛盾真实存在、双重保留不可妥协、投入被严格约束——三者在行业背景（统建系统、公开改造路径验证）下均成立 \\
需求 & 15 条约束 + 4 条操作纪律齐备，每条在方案中有对应条款（约束链 + 切换纪律）——无悬空需求、无过度承诺 \\
技术 & 唯一可行路径（旁路复制）已有行业标准明文背书（NIST SP 800--82r3 / GB/T 20945），难度显式单列工作量，零影响可验收 \\
价值 & 全域感知、多维联动（高频·动态·差异化·信创）全部对应定制计价项——\textbf{定制化流水计算为最优价值}（以前想要也要不到的）；投入产出跨数量级占优、资产可复用 \\
\bottomrule
\end{tabular}
\end{table}

\begin{center}
\large\textbf{A11 动不得 ⟹ 只能旁路拿数据 ⟹ 拿到了 A11 拿不到的数据 ⟹ 做 A11 做不了的应用。}
\end{center}

\vspace{1em}
\noindent\hrulefill

\begin{center}
\small 本论证随项目申报材料提交。与《软件定制开发报价书》、附件②--④、附件⑥、论证报告（第三方证据链）配套使用。\\
编制单位：迪格(杭州)物联科技有限公司 \quad 2026-08
\end{center}

\end{document}
